%%
%% Technical Paper Template for CCS
%%
%% Adapted from SIGCONF Proceedings template
%%

\documentclass[sigconf, pbalance, nonacm]{acmart}
\usepackage{tikz}
\usepackage{listings}
\usepackage{xcolor} 
\usepackage{tabularx}
\usepackage{booktabs}

\AtBeginDocument{%
  \providecommand\BibTeX{{%
    \normalfont{B}\kern-0.5em{\scshape i\kern-0.25em b}\kern-0.8em\TeX}}}

\settopmatter{printacmref=false}
\renewcommand\footnotetextcopyrightpermission{}
\pagestyle{plain}

%%
%% To avoid column breaking
\lstset{
  basicstyle=\ttfamily\small, 
  breaklines=true,            
  breakatwhitespace=true,     
  columns=fullflexible,       
  keepspaces=true,             
  showstringspaces=false,     
  frame=single,               
  commentstyle=\color{gray},   
  keywordstyle=\color{blue}  
}

%% %% %% %% %%
%% Start of the actual paper.
\begin{document}

%%
%% Paper title.
\title{Isolation levels for Sakila Dataset}

%%
%% Authors
\author{Airon Matthew Bantillo}
\email{airon\_bantillo@dlsu.edu.ph}
\affiliation{%
  \institution{De La Salle University}
  \city{Manila}
  \country{Philippines}
}

\author{Bullet Andre Fernandez}
\email{bullet\_fernandez@dlsu.edu.ph}
\affiliation{%
  \institution{De La Salle University}
  \city{Manila}
  \country{Philippines}
}

\author{Charles Sebastian V. Infante}
\email{charles\_infante@dlsu.edu.ph}
\affiliation{%
  \institution{De La Salle University}
  \city{Manila}
  \country{Philippines}
}

\author{Alec Marx Gabriel B. Nono}
\email{alec\_nono@dlsu.edu.ph}
\affiliation{%
  \institution{De La Salle University}
  \city{Manila}
  \country{Philippines}
}

\author{Jean Rondel Ponce}
\email{jean\_ponce@dlsu.edu.ph}
\affiliation{%
  \institution{De La Salle University}
  \city{Manila}
  \country{Philippines}
}

\renewcommand{\shortauthors}{Fernandez et al.}

%%
%% This command processes the author and affiliation and title information and builds the first part of the formatted document.
\maketitle
%%
%%
%% --- Introduction ---
%%
\section{Introduction}

Modern relational database management systems (RDBMS) rely on concurrency controls to enable the simultaneous execution of multiple transactions while ensuring data integrity. ANSI's SQL-92 standards define four transaction isolation levels to balance data consistency and control a system's susceptibility to the following three read anomalies: \textbf{Dirty Reads}, \textbf{Non-Repeatable Reads}, and \textbf{Phantom Reads}. 

This study evaluates the operation behavior of MySQL's \textbf{InnoDB} storage engine under concurrent workload execution. A multi-session testing environment was constructed to observe Multi-Version Concurrency Control behavior across all four isolation levels: \textbf{READ UNCOMMITTED}, \textbf{READ COMMITTED}, \textbf{REPEATABLE READ}, and \textbf{SERIALIZABLE}. 

\subsection {Relational Schema Design}
The experimental domain models MySQL's Sakila database. The relational tables are defined with explicit foreign key constraints and transactional support powered by the \textbf{InnoDB} engine. 

The schema definition for the tracking tables is implemented as follows:

\begin{lstlisting}[language=SQL, caption={DDL Schema Definition for Accounts and Transactions}]
CREATE TABLE IF NOT EXISTS accounts (
    account_id INT PRIMARY KEY AUTO_INCREMENT,
    customer_id INT NOT NULL,
    balance DECIMAL(10, 2) NOT NULL DEFAULT 0.00,
    status VARCHAR(20) NOT NULL DEFAULT 'ACTIVE',
    updated_at TIMESTAMP DEFAULT CURRENT_TIMESTAMP ON UPDATE CURRENT_TIMESTAMP,
    FOREIGN KEY (customer_id) REFERENCES customer(customer_id)
) ENGINE=InnoDB;

CREATE TABLE IF NOT EXISTS transactions (
    transaction_id INT PRIMARY KEY AUTO_INCREMENT,
    account_id INT NOT NULL,
    type VARCHAR(20) NOT NULL, -- 'DEPOSIT', 'WITHDRAWAL', 'TRANSFER'
    amount DECIMAL(10, 2) NOT NULL,
    created_at TIMESTAMP DEFAULT CURRENT_TIMESTAMP,
    FOREIGN KEY (account_id) REFERENCES accounts(account_id)
) ENGINE=InnoDB;
\end{lstlisting} 

\subsection{Initial Seed Data}
The environment was initialized with three core account entities to provide baseline states for concurrent query execution:

\smallskip

\noindent\textbf{Accounts Initial State}
\smallskip

\noindent
\begin{tabularx}{\columnwidth}{XXXX}
101 & 1 & 1000.00 & ACTIVE \\
102 & 2 & 500.00  & ACTIVE \\
103 & 3 & 2500.00 & ACTIVE \\
\end{tabularx}

\vspace{0.4cm} 

\noindent\textbf{Transactions Initial State}
\smallskip

\noindent
\begin{tabularx}{\columnwidth}{XXX}
101 & DEPOSIT & 1000.00 \\
102 & DEPOSIT & 500.00  \\
103 & DEPOSIT & 2500.00 \\
\end{tabularx}
\vspace{0.3cm}

\section{Experiments}

To evaluate concurrency control across the four ANSI SQL-92 isolation levels (\textbf{READ UNCOMMITTED}, \textbf{READ COMMITTED}, \textbf{REPEATABLE READ}, and \textbf{SERIALIZABLE}), a multi-session experimental framework was implemented. The test harness employs two concurrent database connections managed by a multithreaded Python harness with barrier synchronization (\texttt{threading.Barrier}), enforcing precise statement interleaving across dual sessions via controlled \texttt{DO SLEEP()} delays.

A total of 12 baseline experiments ($4\text{ isolation levels} \times 3\text{ read anomalies}$) were systematically executed and logged. Prior to every individual test case, the database state was reset and re-seeded using automated DDL/DML scripts to ensure deterministic initial conditions.

\subsection{Experimental SQL Transactions}
The transactional workloads were designed to isolate the three classical read phenomena. In accordance with SQL coding conventions, transactional commands, query targets, and timing delays are presented below.

\subsubsection{Dirty Read Workload}
A transaction updates Account 101 without committing, while a concurrent transaction reads the uncommitted balance before a rollback occurs.

\smallskip
\noindent\textbf{Real-World Scenario.} Session 1 models a customer withdrawing \$300.00 at an ATM: the balance is debited first, but the machine then encounters a cash-dispensing fault, so the withdrawal is cancelled and rolled back. Session 2 models a bank teller checking the same customer's balance while the withdrawal is still in progress, as when assessing whether a large payment can be approved, and re-checking it moments later. If the teller sees \$700.00, the decision is based on a balance that never officially existed, since the customer's funds remained \$1000.00 throughout.

\begin{lstlisting}[language=SQL, caption={Session 1 (Writer) Transaction for Dirty Read}]
-- Session 1: Writer
SET SESSION TRANSACTION ISOLATION LEVEL <LEVEL>;
START TRANSACTION;
UPDATE accounts SET balance = balance - 300.00 WHERE account_id = 101;
SELECT 'W after update' AS step, balance FROM accounts WHERE account_id = 101;
DO SLEEP(6); -- Hold uncommitted change
ROLLBACK;    -- Reverts update back to 1000.00
SELECT 'W after rollback' AS step, balance FROM accounts WHERE account_id = 101;
\end{lstlisting}

\begin{lstlisting}[language=SQL, caption={Session 2 (Reader) Transaction for Dirty Read}]
-- Session 2: Reader
SET SESSION TRANSACTION ISOLATION LEVEL <LEVEL>;
DO SLEEP(2); -- Wait for Session 1's UPDATE
START TRANSACTION;
SELECT 'R1' AS step, balance FROM accounts WHERE account_id = 101;
DO SLEEP(6); -- Wait for Session 1's ROLLBACK
SELECT 'R2' AS step, balance FROM accounts WHERE account_id = 101;
COMMIT;
\end{lstlisting}

\subsubsection{Non-Repeatable Read Workload}
Session 1 performs an initial read on Account 101, pauses, and re-reads the row. Concurrently, Session 2 modifies and commits the balance during the pause.

\smallskip
\noindent\textbf{Real-World Scenario.} Session 1 models a statement-generation job that reads an account's balance for the summary section and reads it again when producing the closing total of the same document. Session 2 models a \$200.00 payroll deposit that is posted and committed between those two reads. Although the deposit is legitimate, a statement showing \$1000.00 in one section and \$1200.00 in another is internally inconsistent because a single unit of work observed two different committed values.

\begin{lstlisting}[language=SQL, caption={Session 1 (Reader) Transaction for Non-Repeatable Read}]
-- Session 1: Reader
SET SESSION TRANSACTION ISOLATION LEVEL <LEVEL>;
START TRANSACTION;
SELECT 'R1' AS step, balance FROM accounts WHERE account_id = 101;
DO SLEEP(6); -- Pause while Session 2 updates and commits
SELECT 'R2' AS step, balance FROM accounts WHERE account_id = 101;
COMMIT;
\end{lstlisting}

\begin{lstlisting}[language=SQL, caption={Session 2 (Writer) Transaction for Non-Repeatable Read}]
-- Session 2: Writer
SET SESSION TRANSACTION ISOLATION LEVEL <LEVEL>;
DO SLEEP(2); -- Wait for Session 1's first read
START TRANSACTION;
UPDATE accounts SET balance = balance + 200.00 WHERE account_id = 101;
COMMIT;
SELECT 'W committed' AS step, balance FROM accounts WHERE account_id = 101;
\end{lstlisting}

\subsubsection{Phantom Read Workload}
Session 1 executes a range query on accounts with \texttt{balance >= 800.00}, pauses, and re-executes the query. Session 2 concurrently inserts a new qualifying row (Account 104 with balance \$1500.00). In Session 1, an update statement is also executed to evaluate InnoDB current-read behavior under REPEATABLE READ.

\smallskip
\noindent\textbf{Real-World Scenario.} Session 1 models a compliance audit that lists all accounts holding at least \$800.00, re-runs the listing to confirm it, and then flags every qualifying account for review. Session 2 models a branch opening a new account with an initial deposit of \$1500.00 while the audit is running. The new account may appear as an unexpected extra row between the two listings; under REPEATABLE READ, it can remain hidden from both listings yet still be flagged by the update, leaving the audit with a flagged account that it never listed.

\begin{lstlisting}[language=SQL, caption={Session 1 (Reader) Transaction for Phantom Read}]
-- Session 1: Reader
SET SESSION TRANSACTION ISOLATION LEVEL <LEVEL>;
SET SQL_SAFE_UPDATES = 0;
START TRANSACTION;
SELECT 'P1' AS step, account_id, balance FROM accounts WHERE balance >= 800.00 ORDER BY account_id;
DO SLEEP(6); -- Pause while Session 2 inserts row 104
SELECT 'P2' AS step, account_id, balance FROM accounts WHERE balance >= 800.00 ORDER BY account_id;
UPDATE accounts SET status = 'FLAGGED' WHERE balance >= 800.00;
SELECT 'P3' AS step, account_id, balance FROM accounts WHERE balance >= 800.00 ORDER BY account_id;
COMMIT;
\end{lstlisting}

\begin{lstlisting}[language=SQL, caption={Session 2 (Inserter) Transaction for Phantom Read}]
-- Session 2: Inserter
SET SESSION TRANSACTION ISOLATION LEVEL <LEVEL>;
DO SLEEP(2); -- Wait for Session 1's first range read
START TRANSACTION;
INSERT INTO accounts (account_id, customer_id, balance, status) VALUES (104, 4, 1500.00, 'ACTIVE');
COMMIT;
\end{lstlisting}

\subsection{Summary of Baseline Results}
Table~\ref{tab:baseline_matrix} summarizes the empirical findings across all 12 baseline test cases, comparing ANSI-92 expectations with observed InnoDB behavior.

\begin{table*}[t]
\centering
\small
\caption{Comprehensive 12-Case Baseline Concurrency Matrix on MySQL InnoDB}
\label{tab:baseline_matrix}
\begin{tabularx}{\textwidth}{c l l c c c X}
\toprule
\textbf{Case} & \textbf{Anomaly} & \textbf{Isolation Level} & \textbf{ANSI-92} & \textbf{InnoDB Observed} & \textbf{Match} & \textbf{Observed Behavioral Note} \\
\midrule
1 & Dirty Read & READ UNCOMMITTED & Allowed & \textbf{ANOMALY} & Yes & S2 reads uncommitted balance (\$700.00) at $t=2.0$s prior to S1 rollback. \\
2 & Dirty Read & READ COMMITTED & Prevented & None & Yes & S2 reads committed pre-image (\$1000.00) from undo log. \\
3 & Dirty Read & REPEATABLE READ & Prevented & None & Yes & S2 reads from consistent snapshot (\$1000.00). \\
4 & Dirty Read & SERIALIZABLE & Prevented & None & Yes & S2 read is blocked for 4.0s on S1's Exclusive Lock until rollback. \\
\midrule
5 & Non-Repeatable Read & READ UNCOMMITTED & Allowed & \textbf{ANOMALY} & Yes & S1 reads \$1000.00 at $t=0$s and \$1200.00 at $t=6$s ($R_1 \ne R_2$). \\
6 & Non-Repeatable Read & READ COMMITTED & Allowed & \textbf{ANOMALY} & Yes & Second query establishes fresh Read View, reflecting committed \$1200.00. \\
7 & Non-Repeatable Read & REPEATABLE READ & Prevented & None & Yes & Transaction-level Read View preserves \$1000.00 pre-image. \\
8 & Non-Repeatable Read & SERIALIZABLE & Prevented & None & Yes & S1's Shared Lock blocks S2's UPDATE for 4.0s until S1 commits. \\
\midrule
9 & Phantom Read & READ UNCOMMITTED & Allowed & \textbf{ANOMALY} & Yes & Newly inserted Account 104 appears on second range query ($P_2$: 3 rows). \\
10 & Phantom Read & READ COMMITTED & Allowed & \textbf{ANOMALY} & Yes & Fresh statement Read View reveals newly committed Account 104. \\
11 & Phantom Read & REPEATABLE READ & Allowed & \textbf{Write-Only} & Yes & Snapshot hides Account 104 in $P_2$, but range UPDATE touches 104, revealing it in $P_3$. \\
12 & Phantom Read & SERIALIZABLE & Prevented & None & Yes & S1's Next-Key Locks block S2's INSERT for 4.0s until S1 commits. \\
\bottomrule
\end{tabularx}
\end{table*}

\subsection{Two-Column Interleaved Timelines and Empirical Discussion}
The empirical execution results for all 12 baseline configurations are summarized in Table~\ref{tab:baseline_matrix}. To preserve readability in the main text while providing complete empirical transparency, the full two-column interleaved execution traces ($t_x$) across all four isolation levels are grouped into three consolidated configuration tables located in the Appendix:
\begin{itemize}
    \item \textbf{Dirty Read Workload (Cases 1--4):} Documented in Appendix Table~\ref{tab:app_dirty}, illustrating the visibility of uncommitted writes under READ UNCOMMITTED, undo-log pre-image reconstruction under READ COMMITTED and REPEATABLE READ, and Two-Phase Locking contention (4.0s lock block) under SERIALIZABLE.
    \item \textbf{Non-Repeatable Read Workload (Cases 5--8):} Documented in Appendix Table~\ref{tab:app_nonrep}, contrasting statement-level Read Views (fresh snapshot per query under READ COMMITTED) against transaction-level Read View reuse under REPEATABLE READ, and Shared-Lock blocking under SERIALIZABLE.
    \item \textbf{Phantom Read Workload (Cases 9--12):} Documented in Appendix Table~\ref{tab:app_phantom}, demonstrating the classic phantom anomaly in READ UNCOMMITTED and READ COMMITTED, the snapshot hiding and write-only phantom manifestation in REPEATABLE READ, and index gap lock prevention in SERIALIZABLE.
\end{itemize}

\subsubsection{Read View Mechanics: READ COMMITTED vs. REPEATABLE READ}
The difference between READ COMMITTED and REPEATABLE READ for non-repeatable reads stems from allocation policies in Multi-Version Concurrency Control (MVCC) Read View.
Under READ COMMITTED, InnoDB generates statement-level Read Views at the execution of each individual \texttt{SELECT}. 
As session 2 updates to \$1200.00 at $t=2.0$s, Session 1's query at $t=6.0$s generates a 
new Read View that incorporates this committed state, resulting in a non-repeatable read ($R_1 \ne R_2$). 
In contrast, REPEATABLE READ establishes a single, transaction-level Read View upon the execution of a transaction's first read ($t=0.0$s), maintaining it until termination. 
Even if Session 2 commits changes concurrently, Session 1's subsequent read traverses the undo log history
to reconstruct the committed pre-image of \$1000.00, successfully preserving read repeatability.

\subsubsection{The Phantom Phenomenon under REPEATABLE READ}
The ``Write-Only Phantom'' under REPEATABLE READ highlights a distinction between InnoDB's MVCC Snapshot Reads and Current Reads. 
According to ANSI SQL-92's specifications, REPEATABLE READ permits phantom rows on re-queries. 
However, InnoDB mitigates standard read phantoms by utilizing consistent read snapshots; as observed in query $P_2$ at $t=6.0$s, Session 1's plain \texttt{SELECT} 
continues to return only the original two accounts ([101, 103]), completely concealing Session 2's committed insert of Account 104. 
However, when Session 1 subsequently executes a range modification (\texttt{UPDATE accounts SET status = 'FLAGGED' WHERE balance >= 800.00}), the database engine switches to a Current Read to ensure data integrity. 
The \texttt{UPDATE} scans the current clustered index, matches the committed row 104, and updates it. 
Crucially, this operation rewrites row 104's header metadata (\texttt{DB\_TRX\_ID}) with Session 1's own transaction ID. 
Because MVCC rules dictate that a transaction can always observe its own writes, the subsequent query $P_3$ reveals Account 104 alongside accounts 101 and 103. 
The phantom is therefore ``write-only'', as it cannot be triggered through non-locking reads alone but is unmasked once a transaction modifies a phantom record.

\subsubsection{Strict Two-Phase Locking and Empirical Lock Contention}
The 4.0-second execution delay observed in the SERIALIZABLE logs validates Strict Two-Phase Locking (SS2PL) and lock compatibility enforcement. 
In SERIALIZABLE mode, as non-locking MVCC reads are disabled, all \texttt{SELECT} statements acquire Shared Locks (S-locks), while \texttt{UPDATE} statements acquire Exclusive Locks (X-locks). 
In the Dirty Read experiment (Case 4), Session 1 enters its Growing Phase at $t=0.0$s by acquiring an X-lock on Account 101 and holds it across its 6-second sleep. 
When Session 2 issues a read at $t=2.0$s requesting an S-lock, InnoDB blocks the request due to S/X lock conflict. 
Under Strict 2PL, Session 1 only enters its Shrinking Phase upon issuing \texttt{ROLLBACK} at $t=6.0$s, releasing all acquired locks. 
Session 2 is immediately granted the S-lock, logging an exact lock wait duration of 4.0s ($6.0\text{s} - 2.0\text{s}$). 
An identical 4.0s block occurs in Case 8, where Session 1's S-lock blocks Session 2's X-lock update until Session 1 commits.
emonstrating the classic phantom anomaly in READ UNCOMMITTED and READ COMMITTED, the snapshot hiding and write-only phantom manifestation in REPEATABLE READ, and index gap lock prevention in SERIALIZABLE.

\subsection{Explicit Locking Experiments Framework (Member 4)}
To evaluate lock-based concurrency control as an alternative to elevating entire transaction isolation levels, four explicit locking configurations were tested using row-level locks:
\begin{enumerate}
    \item \textbf{Lock Case 1 (Preventing Dirty Reads under READ UNCOMMITTED):} Using \texttt{SELECT ... FOR UPDATE} or table locks.
    \item \textbf{Lock Case 2 (Preventing Non-Repeatable Reads under READ COMMITTED):} Employing \texttt{SELECT ... FOR SHARE} or \texttt{FOR UPDATE} on single-row lookups.
    \item \textbf{Lock Case 3 (Preventing Phantom Reads under READ COMMITTED / REPEATABLE READ):} Employing locking range reads (\texttt{SELECT ... FOR UPDATE WHERE balance >= 800.00}) to invoke explicit Next-Key locking.
    \item \textbf{Lock Case 4 (Performance and Deadlock Benchmarks under Explicit Locking):} Measuring lock-wait latency, lock wait timeouts (\texttt{innodb\_lock\_wait\_timeout}), and throughput trade-offs across concurrent connections.
\end{enumerate}

%% [Member 4 Lock Experiment empirical data, charts, and metrics to be populated here]

%%
%% --- Appendix ---
%%
\appendix
\section{Baseline Concurrency Execution Timelines}
This appendix details the physical query execution traces across dual database sessions (Session 1 and Session 2) captured by the automated concurrency runner. All four ANSI SQL-92 isolation levels (\textbf{READ UNCOMMITTED}, \textbf{READ COMMITTED}, \textbf{REPEATABLE READ}, and \textbf{SERIALIZABLE}) are evaluated across the three concurrency problem configurations. Elapsed execution timestamps ($t_x$) denote seconds from initial barrier synchronization ($t=0.0$s).

\subsection{Dirty Read Baseline Timelines (Cases 1--4)}
Table~\ref{tab:app_dirty} consolidates the execution traces for the Dirty Read workload across all four isolation levels.

\begin{table*}[p]
\centering
\small
\caption{Appendix Table A1: Two-Column Interleaved Timelines for Dirty Read across All Four ANSI Isolation Levels}
\label{tab:app_dirty}
\begin{tabularx}{\textwidth}{p{1.8cm} X X}
\toprule
\textbf{Time ($t_x$)} & \textbf{Session 1: Writer (\texttt{account\_id = 101})} & \textbf{Session 2: Reader (\texttt{account\_id = 101})} \\
\midrule
\multicolumn{3}{l}{\textbf{Configuration 1.1: READ UNCOMMITTED --- Verdict: ANOMALY OBSERVED}} \\
$t = 0.0$s & START TRANSACTION; UPDATE accounts SET balance = balance - 300.00 WHERE account\_id = 101; & DO SLEEP(2); \\
$t = 0.0$s & SELECT 'W after update', balance $\rightarrow$ \textbf{700.00}; DO SLEEP(6); & (waiting on sleep) \\
$t = 2.0$s & (holding uncommitted update) & START TRANSACTION; SELECT 'R1', balance $\rightarrow$ \textbf{700.00} \textbf{[DIRTY READ]}; \\
$t = 6.0$s & ROLLBACK; SELECT 'W after rollback', balance $\rightarrow$ \textbf{1000.00}; & DO SLEEP(6); \\
$t = 8.0$s & (transaction finished) & SELECT 'R2', balance $\rightarrow$ \textbf{1000.00}; COMMIT; \\
\midrule
\multicolumn{3}{l}{\textbf{Configuration 1.2: READ COMMITTED --- Verdict: PREVENTED}} \\
$t = 0.0$s & START TRANSACTION; UPDATE accounts SET balance = balance - 300.00 WHERE account\_id = 101; & DO SLEEP(2); \\
$t = 0.0$s & SELECT 'W after update', balance $\rightarrow$ \textbf{700.00}; DO SLEEP(6); & (waiting on sleep) \\
$t = 2.0$s & (holding uncommitted update) & START TRANSACTION; SELECT 'R1', balance $\rightarrow$ \textbf{1000.00} \textbf{[NO DIRTY READ]}; \\
$t = 6.0$s & ROLLBACK; SELECT 'W after rollback', balance $\rightarrow$ \textbf{1000.00}; & DO SLEEP(6); \\
$t = 8.0$s & (transaction finished) & SELECT 'R2', balance $\rightarrow$ \textbf{1000.00}; COMMIT; \\
\midrule
\multicolumn{3}{l}{\textbf{Configuration 1.3: REPEATABLE READ --- Verdict: PREVENTED}} \\
$t = 0.0$s & START TRANSACTION; UPDATE accounts SET balance = balance - 300.00 WHERE account\_id = 101; & DO SLEEP(2); \\
$t = 0.0$s & SELECT 'W after update', balance $\rightarrow$ \textbf{700.00}; DO SLEEP(6); & (waiting on sleep) \\
$t = 2.0$s & (holding uncommitted update) & START TRANSACTION; SELECT 'R1', balance $\rightarrow$ \textbf{1000.00} \textbf{[NO DIRTY READ]}; \\
$t = 6.0$s & ROLLBACK; SELECT 'W after rollback', balance $\rightarrow$ \textbf{1000.00}; & DO SLEEP(6); \\
$t = 8.0$s & (transaction finished) & SELECT 'R2', balance $\rightarrow$ \textbf{1000.00}; COMMIT; \\
\midrule
\multicolumn{3}{l}{\textbf{Configuration 1.4: SERIALIZABLE --- Verdict: PREVENTED VIA LOCK CONTENTION}} \\
$t = 0.0$s & START TRANSACTION; UPDATE accounts SET balance = balance - 300.00 WHERE account\_id = 101; [Holds X-Lock] & DO SLEEP(2); \\
$t = 0.0$s & SELECT 'W after update', balance $\rightarrow$ \textbf{700.00}; DO SLEEP(6); & (waiting on sleep) \\
$t = 2.0$s & (holding Exclusive X-Lock) & START TRANSACTION; SELECT 'R1', balance $\rightarrow$ \textbf{BLOCKED} (waits for S-Lock); \\
$t = 6.0$s & ROLLBACK; [Releases X-Lock]; SELECT 'W after rollback' $\rightarrow$ \textbf{1000.00}; & \textbf{UNBLOCKED} (waited 4.0s) $\rightarrow$ SELECT 'R1', balance $\rightarrow$ \textbf{1000.00}; DO SLEEP(6); \\
$t = 12.0$s & (transaction finished) & SELECT 'R2', balance $\rightarrow$ \textbf{1000.00}; COMMIT; \\
\bottomrule
\end{tabularx}
\end{table*}

\subsection{Non-Repeatable Read Baseline Timelines (Cases 5--8)}
Table~\ref{tab:app_nonrep} consolidates the execution traces for the Non-Repeatable Read workload across all four isolation levels.

\begin{table*}[p]
\centering
\small
\caption{Appendix Table A2: Two-Column Interleaved Timelines for Non-Repeatable Read across All Four ANSI Isolation Levels}
\label{tab:app_nonrep}
\begin{tabularx}{\textwidth}{p{1.8cm} X X}
\toprule
\textbf{Time ($t_x$)} & \textbf{Session 1: Reader (\texttt{account\_id = 101})} & \textbf{Session 2: Writer (\texttt{account\_id = 101})} \\
\midrule
\multicolumn{3}{l}{\textbf{Configuration 2.1: READ UNCOMMITTED --- Verdict: ANOMALY OBSERVED}} \\
$t = 0.0$s & START TRANSACTION; SELECT 'R1', balance $\rightarrow$ \textbf{1000.00}; DO SLEEP(6); & DO SLEEP(2); \\
$t = 2.0$s & (holding transaction open) & START TRANSACTION; UPDATE accounts SET balance = balance + 200.00 WHERE account\_id = 101; COMMIT; \\
$t = 2.0$s & (holding transaction open) & SELECT 'W committed', balance $\rightarrow$ \textbf{1200.00}; \\
$t = 6.0$s & SELECT 'R2', balance $\rightarrow$ \textbf{1200.00} \textbf{[NON-REPEATABLE READ: $1000 \ne 1200$]}; COMMIT; & (transaction finished) \\
\midrule
\multicolumn{3}{l}{\textbf{Configuration 2.2: READ COMMITTED --- Verdict: ANOMALY OBSERVED}} \\
$t = 0.0$s & START TRANSACTION; SELECT 'R1', balance $\rightarrow$ \textbf{1000.00}; DO SLEEP(6); & DO SLEEP(2); \\
$t = 2.0$s & (holding transaction open) & START TRANSACTION; UPDATE accounts SET balance = balance + 200.00 WHERE account\_id = 101; COMMIT; \\
$t = 2.0$s & (holding transaction open) & SELECT 'W committed', balance $\rightarrow$ \textbf{1200.00}; \\
$t = 6.0$s & SELECT 'R2', balance $\rightarrow$ \textbf{1200.00} \textbf{[NON-REPEATABLE READ: fresh Read View]}; COMMIT; & (transaction finished) \\
\midrule
\multicolumn{3}{l}{\textbf{Configuration 2.3: REPEATABLE READ --- Verdict: PREVENTED}} \\
$t = 0.0$s & START TRANSACTION; SELECT 'R1', balance $\rightarrow$ \textbf{1000.00} [Read View created]; DO SLEEP(6); & DO SLEEP(2); \\
$t = 2.0$s & (holding transaction open) & START TRANSACTION; UPDATE accounts SET balance = balance + 200.00 WHERE account\_id = 101; COMMIT; \\
$t = 2.0$s & (holding transaction open) & SELECT 'W committed', balance $\rightarrow$ \textbf{1200.00}; \\
$t = 6.0$s & SELECT 'R2', balance $\rightarrow$ \textbf{1000.00} \textbf{[REPEATABLE READ: Read View reused]}; COMMIT; & (transaction finished) \\
\midrule
\multicolumn{3}{l}{\textbf{Configuration 2.4: SERIALIZABLE --- Verdict: PREVENTED VIA LOCK CONTENTION}} \\
$t = 0.0$s & START TRANSACTION; SELECT 'R1', balance $\rightarrow$ \textbf{1000.00} [Acquires S-Lock]; DO SLEEP(6); & DO SLEEP(2); \\
$t = 2.0$s & (holding Shared S-Lock) & START TRANSACTION; UPDATE accounts SET balance = balance + 200.00 $\rightarrow$ \textbf{BLOCKED}; \\
$t = 6.0$s & SELECT 'R2', balance $\rightarrow$ \textbf{1000.00}; COMMIT; [Releases S-Lock]; & \textbf{UNBLOCKED} (waited 4.0s) $\rightarrow$ UPDATE completes; COMMIT; \\
$t = 6.0$s & (transaction finished) & SELECT 'W committed', balance $\rightarrow$ \textbf{1200.00}; \\
\bottomrule
\end{tabularx}
\end{table*}

\subsection{Phantom Read Baseline Timelines (Cases 9--12)}
Table~\ref{tab:app_phantom} consolidates the execution traces for the Phantom Read workload across all four isolation levels.

\begin{table*}[p]
\centering
\small
\caption{Appendix Table A3: Two-Column Interleaved Timelines for Phantom Read across All Four ANSI Isolation Levels}
\label{tab:app_phantom}
\begin{tabularx}{\textwidth}{p{1.8cm} X X}
\toprule
\textbf{Time ($t_x$)} & \textbf{Session 1: Reader (\texttt{WHERE balance >= 800.00})} & \textbf{Session 2: Inserter (\texttt{account\_id = 104})} \\
\midrule
\multicolumn{3}{l}{\textbf{Configuration 3.1: READ UNCOMMITTED --- Verdict: ANOMALY OBSERVED}} \\
$t = 0.0$s & START TRANSACTION; SELECT 'P1', account\_id $\rightarrow$ \textbf{[101, 103]} (2 rows); DO SLEEP(6); & DO SLEEP(2); \\
$t = 2.0$s & (holding transaction open) & START TRANSACTION; INSERT INTO accounts VALUES (104, 4, 1500.00, 'ACTIVE'); COMMIT; \\
$t = 6.0$s & SELECT 'P2', account\_id $\rightarrow$ \textbf{[101, 103, 104]} (3 rows) \textbf{[PHANTOM READ]}; & (transaction finished) \\
$t = 6.0$s & UPDATE accounts SET status = 'FLAGGED' WHERE balance >= 800.00; & (transaction finished) \\
$t = 6.0$s & SELECT 'P3', account\_id $\rightarrow$ \textbf{[101, 103, 104]}; COMMIT; & (transaction finished) \\
\midrule
\multicolumn{3}{l}{\textbf{Configuration 3.2: READ COMMITTED --- Verdict: ANOMALY OBSERVED}} \\
$t = 0.0$s & START TRANSACTION; SELECT 'P1', account\_id $\rightarrow$ \textbf{[101, 103]} (2 rows); DO SLEEP(6); & DO SLEEP(2); \\
$t = 2.0$s & (holding transaction open) & START TRANSACTION; INSERT INTO accounts VALUES (104, 4, 1500.00, 'ACTIVE'); COMMIT; \\
$t = 6.0$s & SELECT 'P2', account\_id $\rightarrow$ \textbf{[101, 103, 104]} (3 rows) \textbf{[PHANTOM READ: fresh Read View]}; & (transaction finished) \\
$t = 6.0$s & UPDATE accounts SET status = 'FLAGGED' WHERE balance >= 800.00; & (transaction finished) \\
$t = 6.0$s & SELECT 'P3', account\_id $\rightarrow$ \textbf{[101, 103, 104]}; COMMIT; & (transaction finished) \\
\midrule
\multicolumn{3}{l}{\textbf{Configuration 3.3: REPEATABLE READ --- Verdict: WRITE-ONLY PHANTOM OCCURS}} \\
$t = 0.0$s & START TRANSACTION; SELECT 'P1', account\_id $\rightarrow$ \textbf{[101, 103]} (2 rows) [Snapshot created]; DO SLEEP(6); & DO SLEEP(2); \\
$t = 2.0$s & (holding transaction open) & START TRANSACTION; INSERT INTO accounts VALUES (104, 4, 1500.00, 'ACTIVE'); COMMIT; \\
$t = 6.0$s & SELECT 'P2', account\_id $\rightarrow$ \textbf{[101, 103]} (2 rows) \textbf{[HIDDEN BY MVCC SNAPSHOT]}; & (transaction finished) \\
$t = 6.0$s & UPDATE accounts SET status = 'FLAGGED' WHERE balance >= 800.00; [Current-read updates row 104]; & (transaction finished) \\
$t = 6.0$s & SELECT 'P3', account\_id $\rightarrow$ \textbf{[101, 103, 104]} (3 rows) \textbf{[WRITE-ONLY PHANTOM REVEALED]}; COMMIT; & (transaction finished) \\
\midrule
\multicolumn{3}{l}{\textbf{Configuration 3.4: SERIALIZABLE --- Verdict: PREVENTED VIA NEXT-KEY GAP LOCK}} \\
$t = 0.0$s & START TRANSACTION; SELECT 'P1', account\_id $\rightarrow$ \textbf{[101, 103]} (2 rows) [Holds Next-Key Locks]; DO SLEEP(6); & DO SLEEP(2); \\
$t = 2.0$s & (holding Next-Key Locks on range and gap) & START TRANSACTION; INSERT INTO accounts VALUES (104, 4, 1500.00, 'ACTIVE') $\rightarrow$ \textbf{BLOCKED}; \\
$t = 6.0$s & SELECT 'P2', account\_id $\rightarrow$ \textbf{[101, 103]} (2 rows); & (waiting on Next-Key lock) \\
$t = 6.0$s & UPDATE accounts SET status = 'FLAGGED' WHERE balance >= 800.00; (affected = 2 rows); & (waiting on Next-Key lock) \\
$t = 6.0$s & SELECT 'P3', account\_id $\rightarrow$ \textbf{[101, 103]} (2 rows); COMMIT; [Releases Next-Key Locks]; & \textbf{UNBLOCKED} (waited 4.0s) $\rightarrow$ INSERT completes; COMMIT; \\
\bottomrule
\end{tabularx}
\end{table*}

%%
%% This local block overrides the missing file error and registers your specific citation keys.
\begin{thebibliography}{99}

\bibitem{Silberschatz2020}
Abraham Silberschatz, Henry F. Korth, and S. Sudarshan. 2020.
\newblock {\em Database System Concepts} (7th ed.).
\newblock McGraw-Hill.

\bibitem{RDTArchive2025}
Sakila Database Reference Archive. 2025.
\newblock MySQL Documentation Team.

\bibitem{Selinger1979}
Patricia G. Selinger, et al. 1979.
\newblock Access Path Selection in a Relational Database Management System.
\newblock In {\em Proceedings of the 1979 ACM SIGMOD International Conference on Management of Data}. 23--34.

\end{thebibliography}

\end{document}
